\begin{table}[htbp]
  \centering
  \caption{Impulse Response Function: Effect of $+100$k Troops on $\Delta$ Debt/GDP}
  \label{tab:irf}
  \begin{tabular}{crrrrrr}
    \toprule
    $h$ (years) & IRF & SE & 95\% CI lower & 95\% CI upper & Cum. IRF & Cum. SE \\
    \midrule
    0 & 0.00185 & 0.01252 & -0.02268 & 0.02638 & 0.00185 & 0.01252 \\
    1 & 0.00107 & 0.00722 & -0.01309 & 0.01522 & 0.00292 & 0.01973 \\
    2 & 0.00061 & 0.00417 & -0.00756 & 0.00879 & 0.00353 & 0.02390 \\
    3 & 0.00035 & 0.00241 & -0.00437 & 0.00508 & 0.00388 & 0.02630 \\
    4 & 0.00020 & 0.00139 & -0.00253 & 0.00293 & 0.00409 & 0.02769 \\
    5 & 0.00012 & 0.00081 & -0.00147 & 0.00170 & 0.00420 & 0.02849 \\
    10 & 0.00001 & 0.00005 & -0.00010 & 0.00011 & 0.00435 & 0.02952 \\
    15 & 0.00000 & 0.00000 & -0.00001 & 0.00001 & 0.00436 & 0.02959 \\
    20 & 0.00000 & 0.00000 & -0.00000 & 0.00000 & 0.00436 & 0.02959 \\
    \bottomrule
  \end{tabular}
  \\[2pt] \footnotesize{IRF computed from the MA($\infty$) representation of the ARMA(1,0) part, scaled by $\hat{\beta}$ and the shock size. Shock: a one-time increase of 100,000 troops. Standard errors and confidence intervals via the delta method (numerical Jacobian of the IRF with respect to all estimated parameters). Long-run multiplier (sum over 20 years): 0.00436.}
\end{table}
